\documentclass[a4paper,12pt]{article}
\usepackage{amsmath}
\usepackage{geometry}
\geometry{margin=2.5cm}

\title{Oppsummering: FOPDT-modell og PID-tuning (Lambda-metoden)}
\author{}
\date{}

\begin{document}
\maketitle

\section*{1. FOPDT-modellen}
En First-Order Plus Dead Time (FOPDT)-modell for en trinnrespons er gitt ved:
\[
y(t) =
\begin{cases}
y_{\text{start}}, & t < L \\
y_{\text{start}} + \Delta Y\cdot \big(1 - e^{-\frac{t-L}{T}}\big), & t \ge L
\end{cases}
\]

\noindent hvor:
\begin{itemize}
    \item $y_{\text{start}}$ = startverdi før steget
    \item $\Delta Y = K \cdot \Delta U$ = amplituden (endring i utgang)
    \item $K$ = prosessforsterkning
    \item $\Delta U$ = stegstørrelse i inngangen
    \item $T$ = tidskonstant
    \item $L$ = forsinkelse (dead time)
\end{itemize}

\section*{2. Beregning av parametere}
\begin{align*}
\Delta Y &= y_{\text{slutt}} - y_{\text{start}} \\
K &= \frac{\Delta Y}{\Delta U} \\
\Delta Y &= K \cdot \Delta U
\end{align*}

\noindent Tidskonstanten $\tau$ finnes som:
\[
T = T_{63\%} - \theta
\]
der $t_{63\%}$ er tidspunktet når prosessen har nådd $63\%$ av endringen:
\[
y_{63\%} = y_{\text{start}} + 0.63 \cdot \Delta Y
\]

\section*{3. Lambda-metoden (SIMC tuning)}
For en FOPDT-prosess:
\[
K_p = \frac{\tau}{K(\lambda + L)}, \quad T_i =min(T,\lambda+L), \quad T_d = \frac{L}{2}
\]

\noindent hvor:
\begin{itemize}
    \item $\lambda$ = ønsket lukket sløyfe-tidkonstant (velges av brukeren)
    \item Lav $\lambda$ gir rask respons (aggressiv tuning)
    \item Høy $\lambda$ gir rolig respons (konservativ tuning)
\end{itemize}

\section*{4. Eksempel}
Gitt:
\[
y_{\text{start}} = 0.253, \quad y_{\text{slutt}} = 0.78, \quad \Delta U = 0.1
\]
\[
\Delta Y = 0.527, \quad K = 5.27, \quad L = 3 \text{ s}, \quad T = 2 \text{ s}
\]

\noindent PID-parametere for ulike $\lambda$:
\[
\begin{array}{|c|c|c|c|}
\hline
\lambda & K_p & T_i & T_d \\
\hline
3 & 0.063 & 2 & 1.5 \\
10 & 0.029 & 2 & 1.5 \\
20 & 0.017 & 2 & 1.5 \\
\hline
\end{array}
\]

\end{document}